% Part: many-valued-logic
% Chapter: three-valued-logics
% Section: lukasiewicz

\documentclass[../../../include/open-logic-section]{subfiles}

\begin{document}

\olfileid{mvl}{inf}{luk}

\olsection{د لوکاشېوېچ منطق}

\begin{editorial}
  دا د نامتناهي ارزښته لوکاشېوېچ د منطق د برخې يوه لنډه لومړنۍ مسوده ده۔
\end{editorial}

\begin{defn}\ollabel{def:lukasiewicz} نامتناهي ارزښته د لوکاشېوېچ
منطق~$\LogLuk[\infty]$ په دې ماتريس تعريفېږي:
\begin{enumerate}
  \item معياري بياني ژبه~$\Lang L_0$ چې
  $\lnot$, $\land$, $\lor$, $\lif$ لري۔
  \item د رښتياوالي د قيمتونو سټ~$V_\infty$۔
  \item يوازې~$1$ ټاکل شوے قيمت دے؛ يعنې~$V^+ = \{1\}$۔
  \item صدق تابعې په لاندې تابعو ټاکل کېږي:
  \begin{align*}
    \tf{\lnot}[\LogLuk](x) & = 1 - x\\
    \tf{\land}[\LogLuk](x,y) & = \min(x,y)\\
    \tf{\lor}[\LogLuk](x,y) & = \max(x,y)\\
    \tf{\lif}[\LogLuk](x,y) & = \min(1,1-(x-y)) = \begin{cases}
      1 & \text{که } x \le y\\
      1-(x-y) & \text{په بل حالت کښې.}
    \end{cases}
    \end{align*}
\end{enumerate}
د لوکاشېوېچ $m$-ارزښته منطق همداسې تعريفېږي، خو~$V = V_m$ لري۔
\end{defn}

\begin{prop}
  د~\olref[thr][luk]{def:lukasiewicz} له مخې تعريف شوے منطق~$\LogLuk[3]$
  د~\olref{def:lukasiewicz} له مخې له تعريف شوي~$\LogLuk[3]$ سره يو شان دے۔
\end{prop}

\begin{proof}
  دا د~\olref[thr][luk]{def:lukasiewicz} د نښلوونکو صدق جدولونه
  د~\olref{def:lukasiewicz} له معادلو څخه له ترلاسه شوو جدولونو سره
  په پرتله کولو ليدل کېدے شي:
  \begin{center}
    \begin{tabular}{c|c} 
      $\tf{\lnot}$ & \\ 
      \hline  
      $1$ & $0$ \\ 
      $1/2$ & $1/2$ \\
      $0$ & $1$ 
    \end{tabular}
    \quad
    \begin{tabular}{c|ccc} 
      $\tf{\land}[\LogLuk[3]]$ & $1$ & $1/2$ & $0$ \\ 
      \hline 
      $1$ & $1$ & $1/2$ & $0$ \\ 
      $1/2$ & $1/2$ & $1/2$ & $0$\\ 
      $0$ & $0$ & $0$ & $0$ 
    \end{tabular}
    \\[2ex]
    \begin{tabular}{c|ccc} 
      $\tf{\lor}[\LogLuk[3]]$ & $1$ & $1/2$ & $0$ \\ 
      \hline 
      $1$ & $1$ & $1$ & $1$ \\ 
      $1/2$ & $1$ & $1/2$ & $1/2$ \\
      $0$ & $1$ & $1/2$ & $0$ 
    \end{tabular}
    \quad
    \begin{tabular}{c|ccc} 
      $\tf{\lif}[\LogLuk[3]]$ & $1$ & $1/2$ & $0$ \\ 
      \hline 
      $1$ & $1$ & $1/2$ & $0$ \\ 
      $1/2$ & $1$ & $1$ & $1/2$  \\ 
      $0$ & $1$ & $1$ & $1$ 
    \end{tabular}
  \end{center} 
\end{proof}

\begin{prop}\ollabel{prop:luk-infty-m}
  که~$\Gamma \Entails[\LogLuk[\infty]] !B$ وي، نو د هر~$m \ge 2$
  دپاره~$\Gamma \Entails[\LogLuk[m]] !B$ هم وي۔
\end{prop}

\begin{proof}
  تمرين۔
\end{proof}

\begin{prob}
  \olref[mvl][inf][luk]{prop:luk-infty-m} ثابت کړئ۔
\end{prob}

په حقيقت کښې برعکس لورے هم صدق کوي۔

نامتناهي ارزښته د لوکاشېوېچ منطق تر ټولو مشهور فازي منطق دے۔
د فازي منطق په ليکنو کښې شرطي نښلوونکے ډېر ځله په~$\lnot !A \lor !B$
تعريفېږي۔ په دې سره به يو نامتناهي ارزښته قوي کليني منطق ترلاسه شي۔

\begin{prob}
  وښيئ چې~$(p \lif q) \lor (q \lif p)$ د~$\LogLuk[\infty]$
  تاوتولوژي ده۔
\end{prob}

\end{document}
